% TOP_PROBLEM: {"id": 2582, "title": "Kourovka Notebook Problem 21.73", "category_id": 17, "category": {"id": 17, "name": "group_theory", "display_name": "Group Theory", "description": "Problems about groups, group actions, representations, and related algebraic structures.", "slug": "group-theory", "order_index": 17, "created_at": "2026-07-31T15:26:25.671Z"}, "status": "open", "problem_number": "KOU-21.73", "set_id": 10}
% FIRST_POSTED: 2026-10-01
% ENTRY_KIND: statement_only
% UnsolvedMath ID 2582; source catalogue code KOU-21.73
% !TeX program = lualatex
% Definition and sources only; this file is not an LLM solution attempt.
\documentclass[11pt,a4paper]{article}
\usepackage[margin=25mm]{geometry}
\usepackage{fontspec}
\setmainfont{lmroman10-regular.otf}[BoldFont=lmroman10-bold.otf,
 ItalicFont=lmroman10-italic.otf,BoldItalicFont=lmroman10-bolditalic.otf]
\usepackage{amsmath,amssymb,mathtools}
\usepackage{unicode-math}
\setmathfont{latinmodern-math.otf}
\AtBeginDocument{\let\setminus\smallsetminus}
\usepackage{xurl,hyperref}
\hypersetup{colorlinks=true,urlcolor=blue}
\setcounter{secnumdepth}{0}
\setlength{\parindent}{0pt}
\setlength{\parskip}{5pt}
\setlength{\emergencystretch}{3em}
\begin{document}
\section*{Kourovka Notebook Problem 21.73}\label{2582}
{\small UnsolvedMath ID 2582 · KOU-21.73 · Group Theory · Kourovka Notebook - New Problems, Issue 21}
\subsection{Definitions and mathematical statement}
For integers $m\ge1$ and $0\le r<m$, put
$r(m)=\{r+km:k\in\mathbb Z\}$. Given disjoint residue classes
$r_1(m_1)$ and $r_2(m_2)$, define a permutation $\tau$ of $\mathbb Z$
by interchanging $r_1+km_1$ with $r_2+km_2$ for every $k\in\mathbb Z$
and fixing all other integers. Each such permutation is an involution,
called a class transposition. Set
\[
 \operatorname{CT}(\mathbb Z)
  =\langle\tau_{r_1(m_1),r_2(m_2)}:
                      r_1(m_1)\cap r_2(m_2)=\varnothing\rangle
  \le\operatorname{Sym}(\mathbb Z).
\]
Products mean composition of permutations. Disjointness can be checked
arithmetically: the two residues must differ modulo
$\gcd(m_1,m_2)$.

An input element is encoded by a finite word in class transpositions,
with each letter specified by its four integer parameters. Different
words may represent the same permutation. The conjugacy decision
problem takes two such words, representing $g,h$, and asks whether
\[
 \exists u\in\operatorname{CT}(\mathbb Z)\quad u^{-1}gu=h.
\]
Is there an algorithm that halts and gives the correct answer for
every pair of valid input words?

The possible conjugator may use class transpositions and moduli that
do not occur in either input; it is not restricted to the subgroup
generated by the displayed letters. Conversely, conjugacy in the full
symmetric group of $\mathbb Z$ is insufficient: the conjugator must
belong to $\operatorname{CT}(\mathbb Z)$. No bound on running time or
on the length of a conjugating word is requested.
\subsection{Short English statement}
Is conjugacy decidable for permutations of the integers expressed as finite products of class transpositions?
\subsection{Sources}
\begin{itemize}
\item[S1] ULAM.AI and the UnsolvedMath Contributors, supplied problems.json, record 2582 (KOU-21.73), Kourovka Notebook - New Problems, Issue 21. Catalogue identity and source transcription; CC BY 4.0.
\url{https://huggingface.co/datasets/ulamai/UnsolvedMath}.
\item[S2] The Kourovka Notebook, 21st edition, new problems, Problem 21.73. Expanded formulation of the supplied catalogue statement.
\url{https://arxiv.org/pdf/1401.0300v47}.
\end{itemize}
\end{document}
